% Options for packages loaded elsewhere
% Options for packages loaded elsewhere
\PassOptionsToPackage{unicode}{hyperref}
\PassOptionsToPackage{hyphens}{url}
%
\documentclass[
  ignorenonframetext,
  aspectratio=169,
  russian,
]{beamer}
\newif\ifbibliography
\usepackage{pgfpages}
\setbeamertemplate{caption}[numbered]
\setbeamertemplate{caption label separator}{: }
\setbeamercolor{caption name}{fg=normal text.fg}
\beamertemplatenavigationsymbolshorizontal
% remove section numbering
\setbeamertemplate{part page}{
  \centering
  \begin{beamercolorbox}[sep=16pt,center]{part title}
    \usebeamerfont{part title}\insertpart\par
  \end{beamercolorbox}
}
\setbeamertemplate{section page}{
  \centering
  \begin{beamercolorbox}[sep=12pt,center]{section title}
    \usebeamerfont{section title}\insertsection\par
  \end{beamercolorbox}
}
\setbeamertemplate{subsection page}{
  \centering
  \begin{beamercolorbox}[sep=8pt,center]{subsection title}
    \usebeamerfont{subsection title}\insertsubsection\par
  \end{beamercolorbox}
}
% Prevent slide breaks in the middle of a paragraph
\widowpenalties 1 10000
\raggedbottom
\AtBeginPart{
  \frame{\partpage}
}
\AtBeginSection{
  \ifbibliography
  \else
    \frame{\sectionpage}
  \fi
}
\AtBeginSubsection{
  \frame{\subsectionpage}
}
\usepackage{iftex}
\ifPDFTeX
  \usepackage[T1]{fontenc}
  \usepackage[utf8]{inputenc}
  \usepackage{textcomp} % provide euro and other symbols
\else % if luatex or xetex
  \usepackage{unicode-math} % this also loads fontspec
  \defaultfontfeatures{Scale=MatchLowercase}
  \defaultfontfeatures[\rmfamily]{Ligatures=TeX,Scale=1}
\fi
\usepackage{lmodern}

\ifPDFTeX\else
  % xetex/luatex font selection
\fi
% Use upquote if available, for straight quotes in verbatim environments
\IfFileExists{upquote.sty}{\usepackage{upquote}}{}
\IfFileExists{microtype.sty}{% use microtype if available
  \usepackage[]{microtype}
  \UseMicrotypeSet[protrusion]{basicmath} % disable protrusion for tt fonts
}{}


\usepackage{longtable,booktabs,array}
\usepackage{calc} % for calculating minipage widths
\usepackage{caption}
% Make caption package work with longtable
\makeatletter
\def\fnum@table{\tablename~\thetable}
\makeatother
\usepackage{graphicx}
\makeatletter
\newsavebox\pandoc@box
\newcommand*\pandocbounded[1]{% scales image to fit in text height/width
  \sbox\pandoc@box{#1}%
  \Gscale@div\@tempa{\textheight}{\dimexpr\ht\pandoc@box+\dp\pandoc@box\relax}%
  \Gscale@div\@tempb{\linewidth}{\wd\pandoc@box}%
  \ifdim\@tempb\p@<\@tempa\p@\let\@tempa\@tempb\fi% select the smaller of both
  \ifdim\@tempa\p@<\p@\scalebox{\@tempa}{\usebox\pandoc@box}%
  \else\usebox{\pandoc@box}%
  \fi%
}
% Set default figure placement to htbp
\def\fps@figure{htbp}
\makeatother



\ifLuaTeX
\usepackage[bidi=basic,provide=*]{babel}
\else
\usepackage[bidi=default,provide=*]{babel}
\fi
% get rid of language-specific shorthands (see #6817):
\let\LanguageShortHands\languageshorthands
\def\languageshorthands#1{}


\setlength{\emergencystretch}{3em} % prevent overfull lines

\providecommand{\tightlist}{%
  \setlength{\itemsep}{0pt}\setlength{\parskip}{0pt}}



 

\usepackage[]{csquotes}

\IfFileExists{plex-otf.sty}{
  %% Full TeXlive
  % \usepackage[%
  %   % math,
  %   RM={Scale=0.94},SS={Scale=0.94},SScon={Scale=0.94},TT={Scale=MatchLowercase,FakeStretch=0.9},DefaultFeatures={Ligatures=Common}
  % ]{plex-otf}
}{
  %% TinyTeX
  \usepackage{libertine}
}

%%% Load theme
% https://deic.uab.cat/~iblanes/beamer_gallery/
\IfFileExists{beamerthemegotham.sty}{
  %% Full TeXlive
  \usetheme{gotham}
  \gothamset{
    numbering=totalpagenumber,
    parttocframe default=off,
    sectiontocframe default=off,
    subsectiontocframe default=off,
  }
}{
  %% TinyTeX
  \usetheme{Madrid}
}
\makeatletter
\@ifpackageloaded{caption}{}{\usepackage{caption}}
\AtBeginDocument{%
\ifdefined\contentsname
  \renewcommand*\contentsname{Содержание}
\else
  \newcommand\contentsname{Содержание}
\fi
\ifdefined\listfigurename
  \renewcommand*\listfigurename{Список иллюстраций}
\else
  \newcommand\listfigurename{Список иллюстраций}
\fi
\ifdefined\listtablename
  \renewcommand*\listtablename{Список таблиц}
\else
  \newcommand\listtablename{Список таблиц}
\fi
\ifdefined\figurename
  \renewcommand*\figurename{Рисунок}
\else
  \newcommand\figurename{Рисунок}
\fi
\ifdefined\tablename
  \renewcommand*\tablename{Таблица}
\else
  \newcommand\tablename{Таблица}
\fi
}
\@ifpackageloaded{float}{}{\usepackage{float}}
\floatstyle{ruled}
\@ifundefined{c@chapter}{\newfloat{codelisting}{h}{lop}}{\newfloat{codelisting}{h}{lop}[chapter]}
\floatname{codelisting}{Список}
\newcommand*\listoflistings{\listof{codelisting}{Листинги}}
\makeatother
\makeatletter
\makeatother
\makeatletter
\@ifpackageloaded{caption}{}{\usepackage{caption}}
\@ifpackageloaded{subcaption}{}{\usepackage{subcaption}}
\makeatother

\usepackage{bookmark}
\IfFileExists{xurl.sty}{\usepackage{xurl}}{} % add URL line breaks if available
\urlstyle{same}
\hypersetup{
  pdftitle={Лабораторная работа № 2},
  pdfauthor={Левищева Анастасия Петровна},
  pdflang={ru-RU},
  hidelinks,
  pdfcreator={LaTeX via pandoc}}


\title{Лабораторная работа № 2}
\subtitle{Архитектура компьютеров и операционные системы}
\author{Левищева Анастасия Петровна}
\date{2026-03-07}

\begin{document}
\frame{\titlepage}

\renewcommand*\contentsname{Содержание}
\begin{frame}[allowframebreaks]
  \frametitle{Содержание}
  \setcounter{tocdepth}{1}
  \tableofcontents
\end{frame}
\setcounter{tocdepth}{1}
\tableofcontents
}

\section{Информация}\label{ux438ux43dux444ux43eux440ux43cux430ux446ux438ux44f}

\begin{frame}{Докладчик}
\phantomsection\label{ux434ux43eux43aux43bux430ux434ux447ux438ux43a}
\begin{columns}[c]
\begin{column}{0.7\linewidth}
\begin{itemize}[<+->]
\tightlist
\item
  Левищева Анастасия Петровна
\item
  group NKAbd-02-25
\item
  Российский университет дружбы народов им. П. Лумумбы
\item
  {[}1032252643@rudn.ru{]}
\end{itemize}
\end{column}
\end{columns}
\end{frame}

\section{Цель
работы}\label{ux446ux435ux43bux44c-ux440ux430ux431ux43eux442ux44b}

\begin{frame}{Цель работы}
Целью данной работы является приобретение практических навыков установки
операционной системы на виртуальную машину, настройки минимально
необходимых для дальнейшей работы сервисов.
\end{frame}

\section{Задание}\label{ux437ux430ux434ux430ux43dux438ux435}

\begin{frame}[fragile]{Задание}
\begin{verbatim}
Дождитесь загрузки графического окружения и откройте терминал. В окне терминала проанализируйте последовательность загрузки системы, выполнив команду dmesg. Можно просто просмотреть вывод этой команды:

dmesg | less

Можно использовать поиск с помощью grep:

dmesg | grep -i "то, что ищем"

Получите следующую информацию.
    Версия ядра Linux (Linux version).
    Частота процессора (Detected Mhz processor).
    Модель процессора (CPU0).
    Объём доступной оперативной памяти (Memory available).
    Тип обнаруженного гипервизора (Hypervisor detected).
    Тип файловой системы корневого раздела.
    Последовательность монтирования файловых систем.
\end{verbatim}
\end{frame}

\section{Выполнение лабораторной
работы}\label{ux432ux44bux43fux43eux43bux43dux435ux43dux438ux435-ux43bux430ux431ux43eux440ux430ux442ux43eux440ux43dux43eux439-ux440ux430ux431ux43eux442ux44b}

\begin{frame}{Устанавливаем виртуальную машину.}
\phantomsection\label{ux443ux441ux442ux430ux43dux430ux432ux43bux438ux432ux430ux435ux43c-ux432ux438ux440ux442ux443ux430ux43bux44cux43dux443ux44e-ux43cux430ux448ux438ux43dux443.}
Папка виртуальных машин. Значения по умолчанию Linux: \$HOME/VirtualBox
VMs.
\end{frame}

\begin{frame}{Установочный образ}
\phantomsection\label{ux443ux441ux442ux430ux43dux43eux432ux43eux447ux43dux44bux439-ux43eux431ux440ux430ux437}
Перенесем установочный образ в папку /var/tmp/имя\_пользователя/iso:
mkdir -p \enquote{\$(vboxmanage list systemproperties \textbar{} grep
\enquote*{Default machine folder:} \textbar{} cut -d':' -f2 \textbar{}
tr -d ' ')/iso}

mv Fedora-Sway-Live-x86\_64-39-1.5.iso \enquote{\$(vboxmanage list
systemproperties \textbar{} grep \enquote*{Default machine folder:}
\textbar{} cut -d':' -f2 \textbar{} tr -d ' ')/iso}
\end{frame}

\begin{frame}{Настройка хост-клавиши}
\phantomsection\label{ux43dux430ux441ux442ux440ux43eux439ux43aux430-ux445ux43eux441ux442-ux43aux43bux430ux432ux438ux448ux438}
Хост-клавишей по умолчанию является правый Ctrl.

По умолчанию в дисплейных классах на клавише правый Ctrl находится
переключатель языка ввода.

Эти значения могут конфликтовать.
\end{frame}

\begin{frame}{Графический интерфейс}
\phantomsection\label{ux433ux440ux430ux444ux438ux447ux435ux441ux43aux438ux439-ux438ux43dux442ux435ux440ux444ux435ux439ux441}
В меню выберите Файл, Настройки.

Выберите Ввод, вкладка Виртуальная машина.

Выберите Сочетание клавиш в строке Хост-комбинация.

Нажмите новое сочетание клавиш.

Нажмите ОК, чтобы сохранить изменения.
\end{frame}

\begin{frame}{Командная строка}
\phantomsection\label{ux43aux43eux43cux430ux43dux434ux43dux430ux44f-ux441ux442ux440ux43eux43aux430}
Проверим текущую комбинацию для хост-клавиши:

VBoxManage getextradata global GUI/Input/HostKeyCombination

По умолчанию установлена комбинация 65508, соответствующая правой
клавише Ctrl.

Установите нужную клавишу (в примере клавиша Menu): VBoxManage
setextradata global GUI/Input/HostKeyCombination 65383

\begin{block}{Создадим и запустим виртуальную машину. Запускаем
установочный образ: (рис.~\ref{fig-001}).}
\phantomsection\label{ux441ux43eux437ux434ux430ux434ux438ux43c-ux438-ux437ux430ux43fux443ux441ux442ux438ux43c-ux432ux438ux440ux442ux443ux430ux43bux44cux43dux443ux44e-ux43cux430ux448ux438ux43dux443.-ux437ux430ux43fux443ux441ux43aux430ux435ux43c-ux443ux441ux442ux430ux43dux43eux432ux43eux447ux43dux44bux439-ux43eux431ux440ux430ux437-ux440ux438ux441.-fig-001.}
\begin{figure}

\centering{

\includegraphics[width=0.7\linewidth,height=\textheight,keepaspectratio]{image/1.jpg}

}

\caption{\label{fig-001}Запуск установочного образа}

\end{figure}%
\end{block}

\begin{block}{После установки: (рис.~\ref{fig-002}).}
\phantomsection\label{ux43fux43eux441ux43bux435-ux443ux441ux442ux430ux43dux43eux432ux43aux438-ux440ux438ux441.-fig-002.}
\begin{figure}

\centering{

\includegraphics[width=0.7\linewidth,height=\textheight,keepaspectratio]{image/2.jpg}

}

\caption{\label{fig-002}Вход в ОС под заданной при установке учетной
записью.}

\end{figure}%
\end{block}

\begin{block}{Войдите в ОС под заданной вами при установке учётной
записью.}
\phantomsection\label{ux432ux43eux439ux434ux438ux442ux435-ux432-ux43eux441-ux43fux43eux434-ux437ux430ux434ux430ux43dux43dux43eux439-ux432ux430ux43cux438-ux43fux440ux438-ux443ux441ux442ux430ux43dux43eux432ux43aux435-ux443ux447ux451ux442ux43dux43eux439-ux437ux430ux43fux438ux441ux44cux44e.}
Нажмите комбинацию Win+Enter для запуска терминала.

Переключитесь на роль супер-пользователя: sudo -i
\end{block}
\end{frame}

\begin{frame}{Обновления}
\phantomsection\label{ux43eux431ux43dux43eux432ux43bux435ux43dux438ux44f}
\begin{block}{Установите средства разработки:}
\phantomsection\label{ux443ux441ux442ux430ux43dux43eux432ux438ux442ux435-ux441ux440ux435ux434ux441ux442ux432ux430-ux440ux430ux437ux440ux430ux431ux43eux442ux43aux438}
sudo dnf -y group install development-tools (рис.~\ref{fig-003}).

\begin{figure}

\centering{

\includegraphics[width=0.7\linewidth,height=\textheight,keepaspectratio]{image/3.jpg}

}

\caption{\label{fig-003}Установка средств разработки}

\end{figure}%
\end{block}

\begin{block}{Обновить все пакеты}
\phantomsection\label{ux43eux431ux43dux43eux432ux438ux442ux44c-ux432ux441ux435-ux43fux430ux43aux435ux442ux44b}
sudo dnf -y update
\end{block}
\end{frame}

\begin{frame}{Повышение комфорта работы}
\phantomsection\label{ux43fux43eux432ux44bux448ux435ux43dux438ux435-ux43aux43eux43cux444ux43eux440ux442ux430-ux440ux430ux431ux43eux442ux44b}
Программы для удобства работы в консоли: sudo dnf -y install tmux mc

Другой вариант консоли: sudo dnf -y install kitty
\end{frame}

\begin{frame}{Автоматическое обновление}
\phantomsection\label{ux430ux432ux442ux43eux43cux430ux442ux438ux447ux435ux441ux43aux43eux435-ux43eux431ux43dux43eux432ux43bux435ux43dux438ux435}
При необходимости можно использовать автоматическое обновление.

Установка программного обеспечения: sudo dnf -y install dnf-automatic

Задаёте необходимую конфигурацию в файле /etc/dnf/automatic.conf.

Запустите таймер: sudo systemctl enable --now dnf-automatic.timer
\end{frame}

\begin{frame}{Отключение SELinux}
\phantomsection\label{ux43eux442ux43aux43bux44eux447ux435ux43dux438ux435-selinux}
В данном курсе мы не будем рассматривать работу с системой безопасности
SELinux.

Поэтому отключим его.

\begin{block}{В файле /etc/selinux/config замените значение}
\phantomsection\label{ux432-ux444ux430ux439ux43bux435-etcselinuxconfig-ux437ux430ux43cux435ux43dux438ux442ux435-ux437ux43dux430ux447ux435ux43dux438ux435}
SELINUX=enforcing

на значение

SELINUX=permissive (рис.~\ref{fig-004}).

\begin{figure}

\centering{

\includegraphics[width=0.7\linewidth,height=\textheight,keepaspectratio]{image/4.jpg}

}

\caption{\label{fig-004}Отключение безопасности SELinux}

\end{figure}%
\end{block}

\begin{block}{Перегрузите виртуальную машину:}
\phantomsection\label{ux43fux435ux440ux435ux433ux440ux443ux437ux438ux442ux435-ux432ux438ux440ux442ux443ux430ux43bux44cux43dux443ux44e-ux43cux430ux448ux438ux43dux443}
sudo systemctl reboot
\end{block}
\end{frame}

\begin{frame}{Настройка раскладки клавиатуры}
\phantomsection\label{ux43dux430ux441ux442ux440ux43eux439ux43aux430-ux440ux430ux441ux43aux43bux430ux434ux43aux438-ux43aux43bux430ux432ux438ux430ux442ux443ux440ux44b}
Войдите в ОС под заданной вами при установке учётной записью.

Нажмите комбинацию Win+Enter для запуска терминала.

Запустите терминальный мультиплексор tmux: tmux

Создайте конфигурационный файл
\textasciitilde/.config/sway/config.d/95-system-keyboard-config.conf:
mkdir -p \textasciitilde/.config/sway touch
\textasciitilde/.config/sway/config.d/95-system-keyboard-config.conf

\begin{block}{Отредактируйте конфигурационный файл}
\phantomsection\label{ux43eux442ux440ux435ux434ux430ux43aux442ux438ux440ux443ux439ux442ux435-ux43aux43eux43dux444ux438ux433ux443ux440ux430ux446ux438ux43eux43dux43dux44bux439-ux444ux430ux439ux43b}
\textasciitilde/.config/sway/config.d/95-system-keyboard-config.conf:
exec\_always /usr/libexec/sway-systemd/locale1-xkb-config --oneshot
(рис.~\ref{fig-005}).

\begin{figure}

\centering{

\includegraphics[width=0.7\linewidth,height=\textheight,keepaspectratio]{image/5.jpg}

}

\caption{\label{fig-005}Конфигурационный файл}

\end{figure}%
\end{block}

\begin{block}{Переключитесь на роль супер-пользователя:}
\phantomsection\label{ux43fux435ux440ux435ux43aux43bux44eux447ux438ux442ux435ux441ux44c-ux43dux430-ux440ux43eux43bux44c-ux441ux443ux43fux435ux440-ux43fux43eux43bux44cux437ux43eux432ux430ux442ux435ux43bux44f}
sudo -i

Отредактируйте конфигурационный файл
/etc/X11/xorg.conf.d/00-keyboard.conf: Section \enquote{InputClass}
Identifier \enquote{system-keyboard} MatchIsKeyboard \enquote{on} Option
\enquote{XkbLayout} \enquote{us,ru} Option \enquote{XkbVariant}
\enquote{,winkeys} Option \enquote{XkbOptions}
\enquote{grp:rctrl\_toggle,compose:ralt,terminate:ctrl\_alt\_bksp}
EndSection

Для этого можно использовать файловый менеджер mc и его встроенный
редактор.

Перегрузите виртуальную машину: sudo systemctl reboot
\end{block}
\end{frame}

\begin{frame}{Установка имени пользователя и названия хоста}
\phantomsection\label{ux443ux441ux442ux430ux43dux43eux432ux43aux430-ux438ux43cux435ux43dux438-ux43fux43eux43bux44cux437ux43eux432ux430ux442ux435ux43bux44f-ux438-ux43dux430ux437ux432ux430ux43dux438ux44f-ux445ux43eux441ux442ux430}
Если при установке виртуальной машины вы задали имяпользователя или имя
хоста, не удовлетворяющее соглашению обименовании, то вам необходимо
исправить это.

Запустите виртуальную машину и залогиньтесь.

Нажмите комбинацию Win+Enter для запуска терминала.

Запустите терминальный мультиплексор tmux: tmux

Переключитесь на роль супер-пользователя: sudo -i

Создайте пользователя (вместо username укажите ваш логин в дисплейном
классе): adduser -G wheel username

Задайте пароль для пользователя (вместо username укажите ваш логин в
дисплейном классе): passwd username

Установите имя хоста (вместо username укажите ваш логин в дисплейном
классе): hostnamectl set-hostname username

\begin{block}{Проверьте, что имя хоста установлено верно:}
\phantomsection\label{ux43fux440ux43eux432ux435ux440ux44cux442ux435-ux447ux442ux43e-ux438ux43cux44f-ux445ux43eux441ux442ux430-ux443ux441ux442ux430ux43dux43eux432ux43bux435ux43dux43e-ux432ux435ux440ux43dux43e}
hostnamectl (рис.~\ref{fig-006}).

\begin{figure}

\centering{

\includegraphics[width=0.7\linewidth,height=\textheight,keepaspectratio]{image/6.jpg}

}

\caption{\label{fig-006}Имя хоста}

\end{figure}%
\end{block}
\end{frame}

\begin{frame}{Установка программного обеспечения для создания
документации}
\phantomsection\label{ux443ux441ux442ux430ux43dux43eux432ux43aux430-ux43fux440ux43eux433ux440ux430ux43cux43cux43dux43eux433ux43e-ux43eux431ux435ux441ux43fux435ux447ux435ux43dux438ux44f-ux434ux43bux44f-ux441ux43eux437ux434ux430ux43dux438ux44f-ux434ux43eux43aux443ux43cux435ux43dux442ux430ux446ux438ux438}
Нажмите комбинацию Win+Enter для запуска терминала.

Запустите терминальный мультиплексор tmux: tmux

Переключитесь на роль супер-пользователя: sudo -i
\end{frame}

\begin{frame}{Работа с языком разметки Markdown}
\phantomsection\label{ux440ux430ux431ux43eux442ux430-ux441-ux44fux437ux44bux43aux43eux43c-ux440ux430ux437ux43cux435ux442ux43aux438-markdown}
Средство pandoc для работы с языком разметки Markdown.

Установка с помощью менеджера пакетов: sudo dnf -y install pandoc

Для работы с перекрёстными ссылками мы используем пакет pandoc-crossref.
Пакет pandoc-crossref в стандартном репозитории отсутствует. Придётся
ставить вручную, скачав с сайта
https://github.com/lierdakil/pandoc-crossref.

При установке pandoc-crossref следует обращать внимание, для какой
версии pandoc он скомпилён. Лучше установить pandoc и pandoc-crossref
вручную.

Скачайте необходимую версию pandoc-crossref
(https://github.com/lierdakil/pandoc-crossref/releases). Посмотрите, для
какой версии откомпилён pandoc-crossref. Скачайте соответствующую версию
pandoc (https://github.com/jgm/pandoc/releases). Распакуйте архивы. Обе
программы собраны в виде статически-линкованных бинарных файлов.
Поместите их в каталог /usr/local/bin

\begin{block}{(рис.~\ref{fig-007}).}
\phantomsection\label{ux440ux438ux441.-fig-007.}
\begin{figure}

\centering{

\includegraphics[width=0.7\linewidth,height=\textheight,keepaspectratio]{image/7.jpg}

}

\caption{\label{fig-007}Pandoc и pandoc-crossref}

\end{figure}%
\end{block}
\end{frame}

\begin{frame}{texlive}
\phantomsection\label{texlive}
Установим дистрибутив TeXlive: sudo dnf -y install texlive-scheme-full

\#Домашнее задание

\begin{block}{(рис.~\ref{fig-008}).}
\phantomsection\label{ux440ux438ux441.-fig-008.}
\begin{figure}

\centering{

\includegraphics[width=0.7\linewidth,height=\textheight,keepaspectratio]{image/8.jpg}

}

\caption{\label{fig-008}Домашнее задание(1)}

\end{figure}%
\end{block}

\begin{block}{(рис.~\ref{fig-009}).}
\phantomsection\label{ux440ux438ux441.-fig-009.}
\begin{figure}

\centering{

\includegraphics[width=0.7\linewidth,height=\textheight,keepaspectratio]{image/9.jpg}

}

\caption{\label{fig-009}Домашнее задание(2)}

\end{figure}%
\end{block}
\end{frame}

\section{Выводы}\label{ux432ux44bux432ux43eux434ux44b}

\begin{frame}{Выводы}
В результате выполнения данной лабораторной работы мы провели установку
системы Fedora Sway на виртуальную машину и её первоначальную настройку.
\end{frame}




\end{document}
